\documentclass[11pt,a4paper]{article}

\usepackage[T1]{fontenc}
\usepackage[utf8]{inputenc}
\usepackage[margin=25mm]{geometry}
\usepackage{amsmath,amssymb}
\usepackage{booktabs}
\usepackage{microtype}

\newcommand{\vect}[1]{\mathbf{#1}}
\newcommand{\mat}[1]{\mathbf{#1}}
\newcommand{\vm}{\sigma_{\mathrm{vM}}}
\newcommand{\allow}{\sigma_{\mathrm{allow}}}
\newcommand{\yield}{\sigma_{\mathrm{y}}}

\title{A plane-stress finite element method for\\
       stress-driven Voronoi pocket patterns}
\author{Voronoi Pocket add-in for Autodesk Fusion, version 0.5.0}
\date{November 2026}

\begin{document}
\maketitle

\begin{abstract}
\noindent
This note documents the mathematics behind the stress-driven mode of the
Voronoi Pocket add-in. A flat face of a part is treated as a thin plate
loaded in its own plane. The plate is meshed with constant-strain triangles
on a boundary-fitted equilateral lattice, the linear elastic equations are
solved with a Jacobi-preconditioned conjugate gradient method, and the
von~Mises stress is recovered at the mesh nodes. The stress field then sets
two properties of a Voronoi pocket pattern: the wall thickness of every
cell, from a rule that keeps material in proportion to load, and the cell
size, through a density-weighted centroidal Voronoi relaxation. All formulas
are given as implemented in \texttt{core/stress.py} and
\texttt{core/geometry.py}. The method is validated against closed-form
solutions; its accuracy and its limits are stated with measured numbers.
\end{abstract}

\section{Scope and assumptions}

The calculation is a design aid for deciding where a pocket pattern should
leave more material. It is not a strength verification. The assumptions are:

\begin{enumerate}
  \item \textbf{Plane stress.} The plate is thin, and all forces act in its
        plane, so $\sigma_{zz}=\sigma_{xz}=\sigma_{yz}=0$.
  \item \textbf{Linear elasticity, small displacements,} isotropic material
        with Young's modulus $E$ and Poisson's ratio $\nu$.
  \item \textbf{Solid plate.} The stress is that of the face before any
        pocket is cut. The pockets change the stress field; this is not fed
        back.
  \item \textbf{Static loading.} No dynamics, fatigue or buckling.
\end{enumerate}

\paragraph{Units.} The implementation uses centimetres and newtons, so
stresses are in $\mathrm{N/cm^2}$, with
$1~\mathrm{MPa} = 100~\mathrm{N/cm^2}$.

\section{Governing equations}

Let $\Omega\subset\mathbb{R}^2$ be the face (an outer outline minus holes),
$t$ the plate thickness and $\vect{u}=(u,v)^{\mathsf T}$ the in-plane
displacement. In Voigt notation the strain and stress are
\begin{equation}
  \boldsymbol{\varepsilon} =
  \begin{pmatrix} \varepsilon_{xx}\\ \varepsilon_{yy}\\ \gamma_{xy} \end{pmatrix}
  =
  \begin{pmatrix}
    \partial u/\partial x\\
    \partial v/\partial y\\
    \partial u/\partial y + \partial v/\partial x
  \end{pmatrix},
  \qquad
  \boldsymbol{\sigma} =
  \begin{pmatrix} \sigma_{xx}\\ \sigma_{yy}\\ \sigma_{xy} \end{pmatrix}
  = \mat{D}\,\boldsymbol{\varepsilon},
\end{equation}
with the plane-stress elasticity matrix
\begin{equation}
  \mat{D} = \frac{E}{1-\nu^2}
  \begin{pmatrix}
    1 & \nu & 0\\
    \nu & 1 & 0\\
    0 & 0 & \dfrac{1-\nu}{2}
  \end{pmatrix}
  =
  \begin{pmatrix}
    d_{11} & d_{12} & 0\\
    d_{12} & d_{11} & 0\\
    0 & 0 & d_{33}
  \end{pmatrix}.
  \label{eq:D}
\end{equation}
Equilibrium without body forces is $\nabla\cdot\boldsymbol{\sigma}=\vect{0}$
in $\Omega$. The boundary $\partial\Omega$ is split into a held part
$\Gamma_u$ with $\vect{u}=\vect{0}$, a loaded part $\Gamma_t$ with a given
traction, and a free remainder. The weak form reads: find $\vect{u}$ with
$\vect{u}=\vect{0}$ on $\Gamma_u$ such that for all admissible $\vect{w}$
\begin{equation}
  t\int_\Omega \boldsymbol{\varepsilon}(\vect{w})^{\mathsf T}\,\mat{D}\,
               \boldsymbol{\varepsilon}(\vect{u})\,\mathrm{d}A
  \;=\; \int_{\Gamma_t} \vect{w}^{\mathsf T}\,\vect{q}\,\mathrm{d}s
  \;+\; \sum_k \vect{w}(\vect{x}_k)^{\mathsf T}\,\vect{F}_k ,
  \label{eq:weak}
\end{equation}
where $\vect{q}$ is a line load (force per length) on loaded edges and
$\vect{F}_k$ are forces at points $\vect{x}_k$.

\paragraph{Independence of the stiffness.} With prescribed forces and
homogeneous displacement conditions, scaling $E$ scales the displacements by
$1/E$ and leaves the stresses unchanged. The implementation therefore
solves with a fixed $E$ and takes only $\nu$ from the material. The stress
does depend on $\nu$ wherever the plate is statically indeterminate (near
supports and around holes).

\section{Mesh}

Let $h$ be the target triangle size; the default is the diagonal of the
face's bounding box divided by 65.

\paragraph{Lattice.} Points of an equilateral lattice are laid over the
bounding box:
\begin{equation}
  \vect{p}_{ij} = \Bigl( x_0 + i\,h + \tfrac{h}{2}\,(j \bmod 2),\;
                        y_0 + j\,\tfrac{\sqrt{3}}{2}\,h \Bigr).
\end{equation}
Each lattice cell between rows $j$ and $j+1$ gives two triangles.

\paragraph{Boundary fitting.} A lattice point whose distance to the boundary
is below a reach $r_s = 0.51\,h$ is moved to its nearest boundary point; a
point further away is kept if it lies inside $\Omega$ and discarded
otherwise. The choice $r_s > h/2$ guarantees that of any two neighbouring
lattice points on opposite sides of the boundary at least one is moved, so
the rim of the mesh consists of moved points only and follows the outline.
Moved points closer together than $0.3\,h$ are merged into one. Without
this, the sliver between them would have to be dropped, which leaves a slit
that acts as a crack. Finally, for every outline vertex whose turning angle
exceeds $25^\circ$, the nearest boundary node within $0.75\,h$ is placed
exactly on that vertex, so sharp corners survive.

\paragraph{Filtering.} A triangle is kept if its three nodes are distinct,
its signed doubled area is at least $2\%$ of that of a full lattice triangle
($\tfrac{\sqrt 3}{2}h^2$), and its centroid lies inside $\Omega$. Of the
remaining triangles only the largest edge-connected group is used, because a
fragment attached by a single node could rotate freely and would make the
stiffness matrix singular.

On the test shapes this mesh reproduces the area of the face to within
$0.2\%$, all rim edges lie on the true boundary, and the worst element
quality $q = 4\sqrt{3}\,A/(a^2+b^2+c^2)$ is $0.59$.

\section{Finite element discretisation}

\subsection{Constant-strain triangle}

Each triangle has nodes $1,2,3$ at $(x_i,y_i)$, ordered counter-clockwise,
with linear shape functions~\cite{zienkiewicz}. Define, with cyclic indices,
\begin{equation}
  b_i = y_j - y_k, \qquad c_i = x_k - x_j, \qquad
  2A = x_1 b_1 + x_2 b_2 + x_3 b_3 .
\end{equation}
The strain is constant in the element,
$\boldsymbol{\varepsilon} = \mat{B}\,\vect{u}_e$ with
$\vect{u}_e = (u_1,v_1,u_2,v_2,u_3,v_3)^{\mathsf T}$ and
\begin{equation}
  \mat{B} = \frac{1}{2A}
  \begin{pmatrix}
    b_1 & 0 & b_2 & 0 & b_3 & 0\\
    0 & c_1 & 0 & c_2 & 0 & c_3\\
    c_1 & b_1 & c_2 & b_2 & c_3 & b_3
  \end{pmatrix}.
  \label{eq:B}
\end{equation}
The element stiffness matrix is
$\mat{K}_e = t\,A\,\mat{B}^{\mathsf T}\mat{D}\,\mat{B}$. Written as $2\times2$
blocks coupling node $p$ with node $q$:
\begin{equation}
  \mat{K}_e^{pq} = \frac{t}{4A}
  \begin{pmatrix}
    d_{11}\,b_p b_q + d_{33}\,c_p c_q & d_{12}\,b_p c_q + d_{33}\,c_p b_q\\[2pt]
    d_{12}\,c_p b_q + d_{33}\,b_p c_q & d_{11}\,c_p c_q + d_{33}\,b_p b_q
  \end{pmatrix}.
  \label{eq:Ke}
\end{equation}
The global matrix $\mat{K}$ is assembled by summing these blocks and stored
row by row as sparse lists.

\subsection{Supports}

Every boundary node lying on a held edge has both displacement components
set to zero (optionally only one, for sliding supports). These degrees of
freedom are removed from the system. At least three held degrees of freedom
are required to exclude the three rigid-body motions of the plane.

\subsection{Loads}

\paragraph{Force on an edge.} A total force $\vect{F}$ on an edge is taken
as a uniform line load. Let the edge's boundary nodes, ordered along the
edge, sit at arc lengths $s_1<\dots<s_m$. Node $k$ receives the share of the
edge it covers,
\begin{equation}
  \vect{f}_k = \vect{F}\,\frac{\ell_k}{\sum_j \ell_j},
  \qquad
  \ell_k = \tfrac12\,(s_k - s_{k-1}) + \tfrac12\,(s_{k+1} - s_k),
\end{equation}
with the missing term dropped at the two ends. This is the consistent nodal
load of a uniform traction for linear elements. For a closed ring (a hole
rim) the gap between the last and the first node is split between them. If
several edges are loaded, the force is divided between them in proportion to
their lengths.

\paragraph{Force at a point.} A force on a single node produces a stress
peak that depends only on the mesh. A point force $\vect{F}$ at $\vect{x}_0$
is therefore spread over the nodes within a radius $r_p = 1.5\,h$, with
cone-shaped weights
\begin{equation}
  \vect{f}_k = \vect{F}\,\frac{w_k}{\sum_j w_j},
  \qquad
  w_k = 1 - \frac{\lVert \vect{x}_k - \vect{x}_0\rVert}{r_p}
  \quad\text{for } \lVert \vect{x}_k - \vect{x}_0\rVert < r_p .
\end{equation}
The resultant is exactly $\vect{F}$. This models a force that arrives
through a small pad or bolt head.

\subsection{Solution}

With held degrees of freedom removed, the system
$\mat{K}\,\vect{u} = \vect{f}$ is symmetric positive definite. It is solved
by the conjugate gradient method~\cite{hestenes} with the diagonal of $\mat{K}$ as
preconditioner, $\mat{M} = \operatorname{diag}(\mat{K})$:
\begin{align*}
  &\vect{u}_0 = \vect{0},\quad \vect{r}_0 = \vect{f},\quad
   \vect{z}_0 = \mat{M}^{-1}\vect{r}_0,\quad \vect{p}_0 = \vect{z}_0,\\
  &\alpha_k = \frac{\vect{r}_k^{\mathsf T}\vect{z}_k}
                   {\vect{p}_k^{\mathsf T}\mat{K}\,\vect{p}_k},\qquad
   \vect{u}_{k+1} = \vect{u}_k + \alpha_k\vect{p}_k,\qquad
   \vect{r}_{k+1} = \vect{r}_k - \alpha_k\mat{K}\,\vect{p}_k,\\
  &\vect{z}_{k+1} = \mat{M}^{-1}\vect{r}_{k+1},\qquad
   \beta_k = \frac{\vect{r}_{k+1}^{\mathsf T}\vect{z}_{k+1}}
                  {\vect{r}_k^{\mathsf T}\vect{z}_k},\qquad
   \vect{p}_{k+1} = \vect{z}_{k+1} + \beta_k\vect{p}_k .
\end{align*}
Iteration stops when $\lVert\vect{r}_k\rVert_2 \le 10^{-7}\,\lVert\vect{f}\rVert_2$.
The support reaction is recovered as
$\vect{R} = \sum_{d\in\text{held}} (\mat{K}\vect{u} - \vect{f})_d$ and
serves as a check: it must equal minus the applied force.

\section{Stress recovery}

\subsection{Element stress}

In each triangle the strain follows from \eqref{eq:B},
\begin{equation}
  \varepsilon_{xx} = \frac{1}{2A}\sum_p b_p u_p,\qquad
  \varepsilon_{yy} = \frac{1}{2A}\sum_p c_p v_p,\qquad
  \gamma_{xy} = \frac{1}{2A}\sum_p \bigl(c_p u_p + b_p v_p\bigr),
\end{equation}
and the stress from \eqref{eq:D}:
$\sigma_{xx} = d_{11}\varepsilon_{xx} + d_{12}\varepsilon_{yy}$,
$\sigma_{yy} = d_{12}\varepsilon_{xx} + d_{11}\varepsilon_{yy}$,
$\sigma_{xy} = d_{33}\gamma_{xy}$.

\subsection{Nodal stress}

\paragraph{Interior nodes.} The stress at a node is the area-weighted mean
over the triangles $T$ that share it:
\begin{equation}
  \boldsymbol{\sigma}_n = \frac{\sum_{T\ni n} A_T\,\boldsymbol{\sigma}_T}
                               {\sum_{T\ni n} A_T}.
\end{equation}

\paragraph{Boundary nodes.} A boundary node has triangles on one side only,
so a mean reads too low exactly where stress peaks. For these nodes a plane
is fitted by least squares to the element stresses of a patch $P$ (all
triangles sharing a node with a triangle of $n$), taken at the element
centroids, and evaluated at the node (a simple form of patch
recovery~\cite{zz}). With
$\Delta\vect{x}_T = \vect{x}_T - \vect{x}_n = (\xi_T,\eta_T)$, each stress
component $s$ is modelled as $s \approx a + b\,\xi + c\,\eta$, giving the
normal equations
\begin{equation}
  \begin{pmatrix}
    N & \sum\xi & \sum\eta\\
    \sum\xi & \sum\xi^2 & \sum\xi\eta\\
    \sum\eta & \sum\xi\eta & \sum\eta^2
  \end{pmatrix}
  \begin{pmatrix} a\\ b\\ c \end{pmatrix}
  =
  \begin{pmatrix} \sum s\\ \sum s\,\xi\\ \sum s\,\eta \end{pmatrix},
\end{equation}
with sums over $T\in P$ and $N=|P|$. The nodal value is $a$, obtained by
Cramer's rule. The fit is used when $N\ge5$ and the matrix is not close to
singular; otherwise the area-weighted mean is kept.

\subsection{Equivalent stress and safety factor}

The von~Mises stress in plane stress is
\begin{equation}
  \vm = \sqrt{\sigma_{xx}^2 - \sigma_{xx}\sigma_{yy} + \sigma_{yy}^2
              + 3\,\sigma_{xy}^2}\, .
\end{equation}
Yielding starts where $\vm$ reaches the yield strength $\yield$. The local
safety factor and the allowed stress for a target safety factor $n_t$ are
\begin{equation}
  n(\vect{x}) = \frac{\yield}{\vm(\vect{x})},
  \qquad
  \allow = \frac{\yield}{n_t}.
\end{equation}
Inside a triangle, $\vm$ is interpolated linearly from its three nodal
values with barycentric weights. In the preview, the safety factor is mapped
to a colour position $c = 1 - \tfrac13\log_2 n$, clamped to $[0,1]$, so that
$n\le1$ is red, $n=2$ yellow, $n=4$ green and $n\ge8$ blue.

\section{From stress to pattern}

The pattern consists of the Voronoi cells of $N_c$ seed points, restricted to
the usable area $\Omega_m$ (the face shrunk by the border margin). Each cell
is shrunk by half its wall thickness to give a pocket.

\subsection{Wall thickness}

\paragraph{Rule.} Keep material in proportion to load. Suppose the solid
plate carries $\sigma$ at some place and a fraction $\kappa$ of the material
is kept there. If the remaining ribs carried the load as well as solid
material, their stress would be about $\sigma/\kappa$. Requiring
$\sigma/\kappa\le\allow$ gives the kept fraction
\begin{equation}
  \kappa = \min\!\Bigl(1,\ \frac{\sigma}{\allow}\Bigr).
  \label{eq:kappa}
\end{equation}

\paragraph{Wall for a kept fraction.} Take a cell of area $A_c$ and
characteristic size $s=\sqrt{A_c}$, and walls of thickness $w$, so the cell
loses a strip of $w/2$ along every side. Its pocket then covers about
$(1-w/s)^2$ of the cell, and the kept fraction is
$\kappa = 1-(1-w/s)^2$. Solving for $w$:
\begin{equation}
  w = s\,\bigl(1-\sqrt{1-\kappa}\,\bigr).
  \label{eq:wall}
\end{equation}

\paragraph{As implemented.} For cell $i$, $\sigma_i$ is the largest
interpolated $\vm$ over the centroid and the vertices of the cell clipped to
$\Omega_m$, and
\begin{equation}
  w_i =
  \begin{cases}
    \max\bigl(w_{\min},\ s_i\,(1-\sqrt{1-\kappa_i})\bigr), & \kappa_i<1,\\[3pt]
    2\,s_i \quad\text{(no pocket: the cell stays solid)}, & \kappa_i = 1 .
  \end{cases}
\end{equation}
The value $2s_i$ exceeds the diameter of the largest circle inscribed in any
convex cell of area $s_i^2$, so no pocket remains whatever the cell's shape.
Because every cell gives up half of its own wall, the rib between cells $i$
and $j$ has thickness $\tfrac12(w_i+w_j)$.

\subsection{Cell size}

Cell size follows the load relative to the rest of the part. Let
$\sigma_{5}$ and $\sigma_{95}$ be the 5th and 95th percentiles of the nodal
$\vm$ values; the percentiles keep a single sharp peak from flattening the
scale. With a user-chosen size ratio $R\ge1$ between the widest and the
narrowest cells, define
\begin{equation}
  \beta(\vect{x}) = \operatorname{clamp}\!\Bigl(
      \frac{\vm(\vect{x})-\sigma_{5}}{\sigma_{95}-\sigma_{5}},\,0,\,1\Bigr),
  \qquad
  \omega(\vect{x}) = 1-\Bigl(1-\frac1R\Bigr)\beta(\vect{x}),
  \qquad
  \rho(\vect{x}) = \frac{1}{\omega(\vect{x})^2}.
  \label{eq:density}
\end{equation}
Here $\omega$ is the wanted cell width as a fraction of the widest, and
$\rho\in[1,R^2]$ the wanted number of cells per unit area. If
$\sigma_{95}-\sigma_{5} < 0.1\,\sigma_{95}$ the part is loaded almost evenly
and uniform cells are used.

\paragraph{Seeding.} Seed points are drawn in $\Omega_m$ by rejection
sampling: a uniform candidate $\vect{x}$ is accepted with probability
$\rho(\vect{x})/R^2$.

\paragraph{Relaxation.} The seeds are evened out by Lloyd iterations: each
seed moves to the weighted centroid of its cell,
\begin{equation}
  \vect{x}_i \leftarrow
  \frac{\int_{V_i\cap\Omega_m} \vect{x}\,\phi(\vect{x})\,\mathrm{d}A}
       {\int_{V_i\cap\Omega_m} \phi(\vect{x})\,\mathrm{d}A},
  \qquad \phi = \rho^2 .
\end{equation}
The exponent comes from the asymptotic theory of centroidal Voronoi
tessellations (Gersho's conjecture, see~\cite{du}): in $d$ dimensions the generators of a
tessellation that is centroidal for a weight $\phi$ are distributed with
density proportional to $\phi^{\,d/(d+2)}$. For $d=2$ this is $\phi^{1/2}$, so
weighting by $\rho^2$ keeps the seeds distributed as $\rho$ itself. The
integrals are evaluated by fanning each clipped cell into triangles from its
centroid and sampling $\phi$ at each triangle's centroid.

\section{Validation}

All cases use $\nu = 0.33$. They are part of the automated tests
(\texttt{tests/test\_stress.py}).

\paragraph{Uniform tension.} A strip of height $H$ and thickness $t$ with
sliding supports, pulled by $F$, on a hand-built mesh. The solution
$\sigma_{xx}=F/(Ht)$, $\sigma_{yy}=\sigma_{xy}=0$, elongation $FL/(EHt)$ and
lateral contraction by $\nu$ is reproduced to a relative error below
$10^{-6}$, as expected: linear displacement fields lie in the element's
function space.

\paragraph{Pure shear.} A rectangle sheared by a force $F$ on its top edge of
width $W$: $\sigma_{xy}=F/(Wt)$ within $2\%$ away from the ends.

\paragraph{In-plane bending.} A cantilever of length $L=20$, height
$H=10$, thickness $t=0.5$ (cm) with an end shear force $F=1000$~N. Beam
theory gives the bending stress $\sigma = F(L-x)\,y/I$ with $I=tH^3/12$ and
the neutral-axis shear stress $\tfrac32 F/(Ht)$, hence $\vm=\sqrt3\cdot300 =
519.6~\mathrm{N/cm^2}$ there. The computed values agree within $6\%$ and
$5\%$ respectively on a mesh of about $1500$ nodes.

\paragraph{Plate with a hole.} For a wide plate under remote stress
$\sigma_\infty$ with a circular hole of radius $a$, the Kirsch
solution~\cite{timoshenko} on
the line through the hole centre perpendicular to the load is
\begin{equation}
  \frac{\sigma_{\theta\theta}}{\sigma_\infty}
  = 1 + \frac{a^2}{2r^2} + \frac{3a^4}{2r^4},
\end{equation}
which is $3$ on the rim, $1.65$ at $r=1.4a$ and $1.22$ at $r=2a$. For a
finite width $W$ the gross-section factor on the rim is approximately
$\bigl[2+(1-d/W)^3\bigr]/(1-d/W)$ with $d=2a$ (Heywood, see~\cite{pilkey}), i.e.\ $3.14$ for
$d/W=0.2$. Table~\ref{tab:hole} compares.

\begin{table}[ht]
  \centering
  \begin{tabular}{lrrrr}
    \toprule
    Triangle size $h/a$ & Nodes & On the rim & At $r=1.4a$ & At $r=2a$\\
    \midrule
    0.50 & 966 & 1.85 & 1.65 & 1.29\\
    0.35 & 1951 & 2.21 & -- & --\\
    0.30 & 2606 & 2.38 & 1.58 & 1.25\\
    0.25 & 3742 & 2.60 & -- & --\\
    \midrule
    Theory & & 3.14 & 1.65 & 1.22\\
    \bottomrule
  \end{tabular}
  \caption{Stress in the load direction, divided by the remote stress, for a
           plate of width $10a$ with a central hole. A dash means the value
           was not recorded.}
  \label{tab:hole}
\end{table}

Away from the rim the error is about $5\%$. On the rim the peak is
underestimated by $17$ to $41\%$ and converges slowly, which is the known
behaviour of the constant-strain triangle at a stress concentration. Since
pockets start a border margin away from any edge, the pattern is steered by
the well-resolved part of the field.

\paragraph{Cost.} The implementation is plain Python.
Table~\ref{tab:time} gives the time for one solve.

\begin{table}[ht]
  \centering
  \begin{tabular}{rrr}
    \toprule
    Nodes & CG iterations & Time (s)\\
    \midrule
    812 & 282 & 0.5\\
    1459 & 385 & 1.4\\
    2364 & 455 & 2.5\\
    3634 & 589 & 5.1\\
    5834 & 764 & 11.7\\
    \bottomrule
  \end{tabular}
  \caption{Solve time against mesh size for a $240\times100$~mm plate. The
           iteration count grows roughly like the square root of the number
           of nodes, so the time grows like its power $1.5$.}
  \label{tab:time}
\end{table}

\section{Limitations}

\begin{itemize}
  \item \textbf{No feedback from the pockets.} The safety factor shown is
        that of the solid plate. Verifying the pocketed plate would need a
        mesh fine enough for the thinnest walls, far beyond the mesh sizes
        of Table~\ref{tab:time}.
  \item \textbf{The wall rule is optimistic.} Equation~\eqref{eq:kappa}
        assumes that the ribs carry load as well as solid material. Ribs at
        an angle to the load carry less, and a random two-dimensional cell
        structure deforms partly by bending of its walls. The target safety
        factor has to cover this.
  \item \textbf{Peaks read low} at hole rims and re-entrant corners
        (Table~\ref{tab:hole}), and the stress under a point force depends on
        the spreading radius $r_p$.
  \item \textbf{In-plane only.} Out-of-plane forces, plate bending and the
        floor left under pockets that do not go through are not modelled.
  \item \textbf{Uniform mesh.} There is no local refinement near holes and
        corners.
\end{itemize}

\section*{Symbols}

\begin{tabular}{ll}
  $E$, $\nu$ & Young's modulus, Poisson's ratio\\
  $t$ & plate thickness\\
  $h$ & triangle size of the mesh\\
  $\mat{D}$, $\mat{B}$, $\mat{K}$ & elasticity, strain--displacement and stiffness matrices\\
  $\vm$ & von Mises stress\\
  $\yield$, $\allow$, $n_t$ & yield strength, allowed stress, target safety factor\\
  $\kappa$ & fraction of material kept in a cell\\
  $s$, $w$, $w_{\min}$ & cell size, wall thickness, thinnest wall\\
  $R$ & ratio of the widest to the narrowest cells\\
  $\rho$, $\phi$ & cell density, weight of the Lloyd relaxation\\
\end{tabular}

\begin{thebibliography}{9}
\bibitem{zienkiewicz}
  O.~C. Zienkiewicz, R.~L. Taylor, J.~Z. Zhu,
  \emph{The Finite Element Method: Its Basis and Fundamentals},
  Butterworth-Heinemann.
\bibitem{timoshenko}
  S.~P. Timoshenko, J.~N. Goodier,
  \emph{Theory of Elasticity}, McGraw-Hill.
  (Kirsch solution for a plate with a circular hole.)
\bibitem{pilkey}
  W.~D. Pilkey, D.~F. Pilkey,
  \emph{Peterson's Stress Concentration Factors}, Wiley.
  (Finite-width correction for a hole in a strip.)
\bibitem{hestenes}
  M.~R. Hestenes, E.~Stiefel,
  Methods of conjugate gradients for solving linear systems,
  \emph{J. Res. Natl. Bur. Stand.} 49 (1952).
\bibitem{du}
  Q.~Du, V.~Faber, M.~Gunzburger,
  Centroidal Voronoi tessellations: applications and algorithms,
  \emph{SIAM Review} 41 (1999).
\bibitem{zz}
  O.~C. Zienkiewicz, J.~Z. Zhu,
  The superconvergent patch recovery and a posteriori error estimates,
  \emph{Int. J. Numer. Methods Eng.} 33 (1992).
\end{thebibliography}

\end{document}
